\section{Orchestration}
\label{sec:orchestration}

The orchestrator (\code{python -m training.pipeline}) replaces the v3 shell loop. It
reads one typed configuration file, runs the cycle of \cref{fig:overview-loop}, and keeps
all run state in a single directory. This section covers the cycle and its checkpoints,
the overlap between self-play and training, the replay window, the promotion gate, and
how champions are published. \Cref{sec:infra} covers where the loop runs.

\subsection{The cycle as a checkpointed state machine}
\label{sec:orch-cycle}

A cycle $N$ runs six stages in order:
\begin{enumerate}
  \item \textbf{self-play}: write shard $N$ with the current champion, or with the
        heuristic evaluator before any network exists;
  \item \textbf{train}: start from the champion's weights and train on the replay window
        (\cref{sec:orch-replay});
  \item \textbf{export}: write the candidate ONNX model and check it for parity
        (\cref{sec:learner-export});
  \item \textbf{gate}: play the candidate against the champion (\cref{sec:orch-gate});
  \item \textbf{promote}: on acceptance, copy the candidate to a new immutable generation;
  \item \textbf{publish}: register and release the new champion (\cref{sec:orch-publish}).
\end{enumerate}
If fewer than a configured number of rows exist (\code{min\_rows\_to\_train}), the cycle
records the last five stages as skipped and only accumulates data. The first network is
promoted unconditionally, after a short sanity match against the heuristic, because there
is no champion to gate against.

\subsection{Resumability}
\label{sec:orch-resume}

Cheap hardware gets interrupted: spot instances are reclaimed with a short
warning~\citep{awsspot}, and free CI jobs have hard time limits (\cref{sec:infra-actions}).
v3 restarted the whole cycle after any failure. In v4 a restart resumes at the first
unfinished stage. Five invariants make this safe:
\begin{enumerate}
  \item \emph{Atomic state.} \code{state.json}, which holds the last completed cycle, the
        stages done in the current cycle and the champion lineage, is written to a
        temporary file and renamed into place. Shards (\cref{sec:selfplay-shards}) and
        history records are written the same way.
  \item \emph{Mark after success.} A stage is recorded as done only after it finishes.
        Counters such as the lifetime row count change in the same write, so a crash
        between the work and the write cannot count rows twice.
  \item \emph{Idempotent gating.} Each arena batch writes its own report file. A resumed
        gate reuses finished batches and replays nothing.
  \item \emph{Mid-run training resume.} If a previous training attempt left a checkpoint,
        the retry passes \code{--continue} so the learner restores optimizer and schedule
        state (\cref{sec:learner-opt}).
  \item \emph{Clean termination.} \code{SIGTERM}, sent by \code{docker stop}, preemption
        notices and CI timeouts, is turned into an orderly shutdown that stops background
        self-play and inference servers.
\end{enumerate}
A test injects a failure into the training stage, starts a new orchestrator process on
the same directory, and checks that self-play is not repeated and that rows are counted
once (\cref{sec:pipeline-eval}).

\subsection{Overlapping self-play with training}
\label{sec:orch-overlap}

In \katago{} and in distributed reinforcement-learning systems, data generation runs
continuously and apart from the learner~\citep{wu2019accelerating,espeholt2018impala,horgan2018distributed}.
A single machine with a GPU is in the same situation on a small scale: self-play is mostly
CPU work, training is mostly GPU work. With \code{overlap\_selfplay} enabled, the
orchestrator starts cycle $N{+}1$'s self-play in the background as soon as cycle $N$'s
self-play finishes, using the champion at that moment. The games are therefore at most one
generation stale, and the stage-$N{+}1$ self-play step only waits for whatever is left.
Overlap is off in the CPU-only configurations, because there both workloads compete for the
same cores.

Prefetched games are written to a staging directory and moved into the replay
directory only when their own cycle begins. Without that, the trainer of cycle $N$ could
see cycle $N{+}1$'s shard: the learner reads the shard directory newest-first, so a
prefetch that finished early displaced the very data the cycle was meant to learn from,
and whether it did so depended on wall-clock timing.

\subsection{Replay as a window over immutable shards}
\label{sec:orch-replay}

Experience replay reuses recent data across updates~\citep{lin1992self,mnih2015human}.
v3 kept one JSONL file, capped at 50{,}000 rows, and rewrote the whole file whenever it
trimmed it. v4 never rewrites data. The buffer is the newest shards whose rows fill the
current window, and the learner receives the shard directory plus the window size in rows
(\code{--window-rows}). The window follows \katago{}'s power law, \cref{eq:window}, applied
to the lifetime row count $N_{\mathrm{total}}$. The anchor $c$ is \code{min\_window\_rows}:
20{,}000 on CPU and 100{,}000 on a GPU, against \katago{}'s 250{,}000. The window is capped at
\code{max\_window\_rows}. The orchestrator owns this value because only it knows the lifetime
total, including shards that have left the live directory. Shards older than
\code{archive\_factor} times the window move to \code{selfplay/archive/}, which keeps the
directory, and therefore any CI cache (\cref{sec:infra-actions}), bounded.

\subsection{Sequential gating}
\label{sec:orch-gate}

\paragraph{Model.}
Under the Elo model~\citep{elo1978rating}, a candidate that is $\Delta$ Elo stronger than
the champion scores
\begin{equation}
  p(\Delta) = \frac{1}{1 + 10^{-\Delta/400}}
  \label{eq:elo-score}
\end{equation}
per game. On $9\times9$ with komi $6.5$ a game cannot be drawn, so a gate is a sequential
test on Bernoulli trials. In what follows, $w$ and $\ell$ are the candidate's wins and losses
so far.

\paragraph{Previous gates.}
\katago{} promotes a candidate that wins at least 100 of 200
games~\citep[App.~E]{wu2019accelerating}. This is a non-regression check: its purpose is to
keep a clearly worse network out of self-play, and at $\Delta=0$ it accepts about half the
time. The v3 gate was stricter: a win rate of at least $0.55$ with a Wilson lower
bound~\citep{wilson1927probable} above $0.5$. But the Go arena also \emph{stopped early}:
after 20 games it rejects once $0.55$ is out of reach, and after 100 games it accepts as
soon as both conditions hold. Checking a significance condition after every game inflates
the false-positive rate; this is the repeated-significance-testing problem of
\citet{armitage1969repeated}. \Cref{tab:orch-gate-oc} shows the effect: a v3 candidate with
no real gain is promoted $8.0\%$ of the time, almost three times the $2.8\%$ of a single
test at 200 games with the same bound.

\paragraph{SPRT.}
v4 uses Wald's sequential probability ratio test~\citep{wald1945sequential}, the procedure
Stockfish's Fishtest uses to accept engine changes~\citep{fishtest}. It compares
$H_0\colon \Delta=\Delta_0$ against $H_1\colon \Delta=\Delta_1$ using the log-likelihood
ratio
\begin{equation}
  \Lambda(w,\ell) = w \log\frac{p_1}{p_0} + \ell \log\frac{1-p_1}{1-p_0},
  \qquad p_i = p(\Delta_i),
  \label{eq:llr}
\end{equation}
where draws, if there were any, count as half a win and half a loss. The test accepts
$H_1$ when $\Lambda \ge \log\frac{1-\beta}{\alpha}$ and accepts $H_0$ when
$\Lambda \le \log\frac{\beta}{1-\alpha}$. The defaults are $\Delta_0=0$, $\Delta_1=35$ and
$\alpha=0.05$, $\beta=0.10$, which puts the bounds at $-2.25$ and $+2.89$. Two
practical details:
\begin{itemize}
  \item \emph{Batches.} The test is checked after arena batches of 40 games, not after
        every game. A full batch has equal numbers of games with each colour, and the
        engine can batch its network evaluations. Each batch is one Go arena call, so the arena's in-batch early reject
        still applies. Stopping a batch early does not bias the likelihood ratio, because
        the next batch simply continues the same sequence of games.
  \item \emph{Cap.} The gate plays at most 600 games by default. If the test is still undecided
        there, the v3 rule decides.
\end{itemize}

\paragraph{Exact operating characteristics.}
We computed each gate's acceptance probability $P_{\mathrm{acc}}$ and expected number of
games $E[n]$ exactly, by dynamic programming over (games played, wins). The computation
calls the orchestrator's own decision code and mirrors the Go arena's stopping rule, so it
describes the gates as they run (\code{paper/analysis/gate\_oc.py}). A 20{,}000-trial Monte
Carlo simulation of the same procedures agrees within sampling error. The only modelling
simplification is that games finish in order, whereas the real arena plays them in
parallel. \Cref{tab:orch-gate-oc,fig:orch-gate-oc} show the results.

\begin{table}[htbp]
\centering
\small
\begin{tabular}{r rr rr rr rr}
\toprule
 & \multicolumn{2}{c}{\katago{} 100/200} & \multicolumn{2}{c}{\gofer{} v3} & \multicolumn{2}{c}{v4 SPRT (default)} & \multicolumn{2}{c}{v4 SPRT (Actions)} \\
\cmidrule(lr){2-3}\cmidrule(lr){4-5}\cmidrule(lr){6-7}\cmidrule(lr){8-9}
True $\Delta$Elo & $P_{\mathrm{acc}}$ & $E[n]$ & $P_{\mathrm{acc}}$ & $E[n]$ & $P_{\mathrm{acc}}$ & $E[n]$ & $P_{\mathrm{acc}}$ & $E[n]$ \\
\midrule
$-100$ & 0.000 & 200 & 0.000 & 142 & 0.000 & 94 & 0.000 & 61 \\
$-50$ & 0.025 & 200 & 0.001 & 159 & 0.000 & 150 & 0.000 & 94 \\
$-20$ & 0.228 & 200 & 0.017 & 171 & 0.001 & 254 & 0.003 & 138 \\
$+0$ & 0.528 & 200 & 0.080 & 175 & 0.023 & 397 & 0.029 & 176 \\
$+20$ & 0.812 & 200 & 0.249 & 172 & 0.233 & 492 & 0.149 & 205 \\
$+35$ & 0.933 & 200 & 0.452 & 161 & 0.605 & 444 & 0.349 & 209 \\
$+50$ & 0.982 & 200 & 0.672 & 145 & 0.899 & 339 & 0.605 & 196 \\
$+100$ & 1.000 & 200 & 0.991 & 105 & 1.000 & 154 & 0.993 & 116 \\
$+200$ & 1.000 & 200 & 1.000 & 100 & 1.000 & 85 & 1.000 & 58 \\
\bottomrule
\end{tabular}

\caption{Exact acceptance probability and expected games for four gates as a function of
the candidate's true advantage. v3 includes the Go arena's interim stops. The v4 columns
use SPRT$(0,\Delta_1)$ with $\alpha=\beta=0.05$: the default configuration ($\Delta_1=35$,
batches of 40, cap 600, $\beta=0.10$) and the free-runner configuration of
\cref{sec:infra-actions} ($\Delta_1=50$, batches of 24, cap 240).}
\label{tab:orch-gate-oc}
\end{table}

\begin{figure}[htbp]
\centering
\begin{tikzpicture}
\begin{groupplot}[
  group style={group size=2 by 1, horizontal sep=1.6cm},
  width=0.47\textwidth, height=5.2cm,
  xlabel={True advantage $\Delta$ (Elo)},
  xmin=-150, xmax=250, grid=major, grid style={gray!20},
  tick label style={font=\footnotesize}, label style={font=\small},
  legend style={font=\footnotesize, draw=none, fill=none},
  every axis plot/.append style={thick, mark size=1.4pt},
]
\nextgroupplot[ylabel={$P(\text{accept})$}, ymin=0, ymax=1.02,
  legend to name=gatelegend, legend columns=4,
  legend style={/tikz/every even column/.append style={column sep=0.5cm}}]
\addplot[black, dashed] table[x=elo, y=katagoAcc, col sep=comma] {figures/gate_oc.csv};
\addlegendentry{\katago{} 100/200}
\addplot[orange!85!black, mark=square*, mark repeat=3] table[x=elo, y=vthreeAcc, col sep=comma] {figures/gate_oc.csv};
\addlegendentry{\gofer{} v3}
\addplot[blue!75!black, mark=*, mark repeat=3] table[x=elo, y=vfourAcc, col sep=comma] {figures/gate_oc.csv};
\addlegendentry{v4 SPRT (default)}
\addplot[teal!70!black, dotted, mark=triangle*, mark repeat=3] table[x=elo, y=actionsAcc, col sep=comma] {figures/gate_oc.csv};
\addlegendentry{v4 SPRT (Actions)}
\nextgroupplot[ylabel={$E[\text{games}]$}, ymin=0, ymax=420]
\addplot[black, dashed] table[x=elo, y=katagoN, col sep=comma] {figures/gate_oc.csv};
\addplot[orange!85!black, mark=square*, mark repeat=3] table[x=elo, y=vthreeN, col sep=comma] {figures/gate_oc.csv};
\addplot[blue!75!black, mark=*, mark repeat=3] table[x=elo, y=vfourN, col sep=comma] {figures/gate_oc.csv};
\addplot[teal!70!black, dotted, mark=triangle*, mark repeat=3] table[x=elo, y=actionsN, col sep=comma] {figures/gate_oc.csv};
\end{groupplot}
\end{tikzpicture}

\vspace{2pt}
\ref*{gatelegend}
\caption{Operating characteristics of the four gates in \cref{tab:orch-gate-oc}: the
probability of promotion (left) and the expected gate length (right), against the
candidate's true advantage.}
\label{fig:orch-gate-oc}
\end{figure}

\paragraph{What SPRT buys, and what it costs.}
Three findings stand out.
\begin{itemize}
  \item \emph{Error control.} SPRT keeps the false-promotion rate close to its nominal
        level: $2.4\%$ at $\Delta=0$, against v3's $8.0\%$. It also has more power for
        real gains: at $\Delta=+50$ it promotes $90\%$ of candidates, against $67\%$.
  \item \emph{Not uniformly cheaper.} Clearly worse candidates are rejected sooner than
        in v3 (84 against 142 expected games at $-100$), and very strong ones are
        accepted sooner at $+200$ (84 against 100). But near the indifference region SPRT
        plays more than twice as many games (380 against 175 at $\Delta=0$). The v4 gate is
        an error-control improvement, not a saving in games.
  \item \emph{The cap matters.} With the default cap of 600 games, a candidate that is
        truly $+35$ Elo stronger is promoted $61\%$ of the time, below the nominal
        $1-\beta=0.90$, because many such tests reach the cap undecided. Raising the cap
        is the main lever: at 800 games the same candidate is promoted $74\%$ of the time,
        which is what the GPU configuration uses. The free-runner configuration's cap of
        240 can only detect large gains ($36\%$ at $+35$, $61\%$ at $+50$). On a
        small budget, that configuration can only reliably detect large gains. This is
        acceptable early in a run, when each generation typically gains a lot, but it
        should be revisited once gains become small.
\end{itemize}
\Cref{tab:orch-gate-oc} also shows what \katago{}'s gate is for. It accepts $23\%$ of
candidates that are \emph{20 Elo worse}. That is consistent with its stated purpose: a
cheap guard against large regressions in a run where most candidates improve.

\subsection{Promotion, the Elo ladder and publishing}
\label{sec:orch-publish}

\paragraph{Generations.}
A promoted candidate becomes generation $g$, stored as three immutable files: the ONNX
model, the weights, and the full checkpoint. Nothing overwrites or deletes a
generation. Numbers are never reused, even after a rollback. Each generation carries a
\emph{ladder} Elo: the previous champion's value plus the gain measured at its gate, with
generation~1 at~0. The ladder is only a progress indicator. It is biased upwards because a
candidate is recorded only when its measured gain was large enough to pass the gate. It is
also blind to intransitivity. A proper rating would fit all cross-generation games jointly,
as \katago{} did with BayesElo~\citep{coulom2010bayeselo,wu2019accelerating}; we leave this
for future work.

\paragraph{Registry and releases.}
The publish stage copies each new champion into a registry directory with a JSON card
(gate statistics, ladder Elo, SHA-256) and appends it to an index. The index tracks the
current \code{best} and the \code{previous\_best}. An optional \emph{alias}, a fixed
path such as \code{models/gofer-9x9-best.onnx} that play and analysis commands can point
at, is overwritten with each new best. The registry is where the history is kept. Optionally, each champion is also published as its
own GitHub Release, so older champions remain downloadable. A failed upload is recorded and
retried later; it never stops training.

\paragraph{Guarding against a false best.}
A gate can pass by chance ($2.4\%$ at $\Delta=0$ in \cref{tab:orch-gate-oc}), or a
candidate can beat the current champion while losing to earlier ones. Before publishing, an
optional \emph{regression match} plays the new champion against the champion
\code{regression\_lookback} generations earlier. If the new champion scores below a
threshold, it is recorded as \emph{held}: the alias and the latest release do not move. It
can also be demoted, so that self-play falls back to the previous champion. A
\code{rollback} command re-points \code{best}, the alias and the latest release to any
earlier generation, and can optionally make that generation the training champion again.
Tests cover all of these paths (\cref{sec:pipeline-eval}).
